%%%%%%%%%%%%%%%%%%%%%%%%%%%%%%%%%%%%%%%%%
% Их сургуулийн оюутны тезис  
% LaTeX Загвар
% Version 2.3 (25/3/16)
%
% Энэ загвар нь дараах сайтаас авсан загварын монгол хувилбар юм.
% http://www.LaTeXTemplates.com
%
% Version 2.x major modifications by:
% Vel (vel@latextemplates.com)
%
% Анхдагч загварын эх үүсвэр:
% Steve Gunn (http://users.ecs.soton.ac.uk/srg/softwaretools/document/templates/)
% Sunil Patel (http://www.sunilpatel.co.uk/thesis-template/)
%
% Загварын лиценз:
% CC BY-NC-SA 3.0 (http://creativecommons.org/licenses/by-nc-sa/3.0/)
%
%%%%%%%%%%%%%%%%%%%%%%%%%%%%%%%%%%%%%%%%%

%-------------------------------------------------------------------------------
%	PACKAGES AND OTHER DOCUMENT CONFIGURATIONS
%-------------------------------------------------------------------------------

\documentclass[
12pt, % Баримтын фонтын хэмжээ, сонголт: 10pt, 11pt, 12pt
oneside, % Хоёр талаар хэвлэж үдэхээр тохируулсан. Нэг тал бол комментыг арилга
%chapterinoneline,% Нэг мөрөнд бүлгийн дугаар, нэрийг гаргах
english, % babel багцын хэлний тохиргоо
onehalfspacing, % Мөр хоорондын зай. Сонголтууд: singlespacing, onehalfspacing, doublespacing
%draft, % Ноорог горимд шилжихийн тулд комментыг арилга(зураг, холбоос, hboxes гарахгүй)
nolistspacing, % Хэрэв мөр хоорондын зай onehalfspacing эсвэл doublespacing бол, жагсаалтын мөр хоорондын зайг single болгохын тулд комментыг арилга
%liststotoc, % Зураг/хүснэгт/бусад жагсаалтыг гарчигт оруулахын тулд комментыг арилга
%toctotoc, % Uncomment to add the main table of contents to the table of contents
%parskip, % Параграф хооронд зай оруулахын тулд комментыг арилга
%nohyperref, % hyperref багцыг ачаалахгүй бол комментыг арилга
headsepline, % Толгой мөрийн доогуур шугам татахын тулд комментыг арилга
]{MUST-Thesis} % Энэ класс файл нь баримтын бүтцийг тодорхойлно

\usepackage[utf8]{inputenc} % Олон улсын тэмдэгт оруулахад хэрэгтэй
\usepackage[T2A]{fontenc} % Олон улсын тэмдэгтийн гаралтын кодчилол
\usepackage[mongolian]{babel}

\usepackage{titlesec}
\usepackage[table]{xcolor}
\usepackage{rotating} % эргүүлэх
\usepackage{enumitem}
\setlist{noitemsep}

\let\oldtabular\tabular
\let\endoldtabular\endtabular
\renewenvironment{tabular}
   {\bgroup\singlespacing\oldtabular}%
   {\endoldtabular\egroup}
\let\oldverse\verse
\let\endoldverse\endverse
\renewenvironment{verse}
   {\bgroup\singlespacing\oldverse}%
   {\endoldverse\egroup}

\usepackage{caption,longtable,ltcaption}
\DeclareCaptionJustification{nohyphen}{\hyphenpenalty=10000}
\captionsetup{justification=nohyphen}

\usepackage{pgf}
\usepackage{tikz} % зураг зурах
\usetikzlibrary{shapes,arrows,automata}

\definecolor{OliveGreen}{cmyk}{0.64,0,0.95,0.40}
\definecolor{CadetBlue}{cmyk}{0.62,0.57,0.23,0}
\definecolor{Cover}{RGB}{102,255,255}
\definecolor{Chapter}{RGB}{255,220,180}

\definecolor{lightlightgray}{gray}{0.9}
\usepackage{listings}
\renewcommand{\lstlistingname}{Эх код} % Програмын эх код хэвлэх
\lstset{
language=C,                             % Code langugage
basicstyle=\linespread{1.0}\ttfamily,                   % Code font
keywordstyle=\color{OliveGreen},        % Keywords font ('*' = uppercase)
commentstyle=\color{CadetBlue},         % Comments font
numbers=left,                           % Line nums position
numberstyle=\tiny,                      % Line-numbers fonts
stepnumber=1,                           % Step between two line-numbers
numbersep=5pt,                          % How far are line-numbers from code
backgroundcolor=\color{lightlightgray}, % Choose background color
frame=none,                             % A frame around the code
tabsize=2,                              % Default tab size
captionpos=t,                           % Caption-position = bottom
breaklines=true,                        % Automatic line breaking?
breakatwhitespace=false,                % Automatic breaks only at whitespace?
showspaces=false,                       % Dont make spaces visible
showtabs=false,                         % Dont make tabls visible
columns=flexible,                       % Column format
morekeywords={__global__, __device__},  % CUDA specific keywords
}

\usepackage[autostyle=false]{csquotes} % Ном зүйд хэлнээс хамаарсан хашилт оруулахад хэрэгтэй

\usepackage[backend=bibtex,natbib=true,sorting=none,sortcites]{biblatex} % Ном зүйд bibtex -г ашиглах

\addbibresource{references.bib} % Ном зүйн файл

\newcommand{\authorshipname}{Зохиогч эрхийн хамгаалал}
\newcommand{\abbrevname}{Товчилсон үгс}
\newcommand{\constantsname}{Физик тогтмолууд}
\newcommand{\symbolsname}{Таних тэмдэгтүүд}
\newcommand{\acknowledgementname}{Талархал}
\newcommand{\abstractname}{Хураангуй}
 
%-------------------------------------------------------------------------------
%	THESIS INFORMATION
%-------------------------------------------------------------------------------

\thesistitle{Багшийн журнал вэб системийн хөгжүүлэлт} % Таны ажлын нэр, нүүр болон хураангуй хуудсанд ашигласан. Өөр газарт бол \ttitle командыг хэрэглэнэ
\thesistype{Бакалаврын төгсөлтийн ажил} % Удирдагчийн нэр, нүүр хуудсанд ашиглана. Дурын газарт бол \thesisname командыг хэрэглэнэ
\supervisor{Магистр Б.Мөнхбуян} % Удирдагчийн нэр, нүүр хуудсанд ашиглана. Дурын газарт бол \supname командыг хэрэглэнэ
\reader{Проф. PhD. А.Батмөнх} % Шүүмжлэгчийн нэр, Дурын газарт бол \readname командыг хэрэглэнэ
\advisor{Доктор (Ph.D) Ц.Ганбат, Магистр Э.Баатархүү} % Зөвлөгчийн нэр, Дурын газарт бол \advicename командыг хэрэглэнэ
\degreeind{D523910} % Мэргэжлийн индекс, нүүр болон хураангуй хуудсанд ашигласан. Өөр газарт бол \degreeid командыг хэрэглэнэ
\degree{Мэдээллийн систем} % Боловсролын зэрэг, нүүр болон хураангуй хуудсанд ашигласан. Өөр газарт бол \degreename командыг хэрэглэнэ
\authorshort{Г.Болд} % Таны товч нэр, нүүр болон хураангуй хуудсанд ашигласан. Дурын газарт бол \shortname командыг хэрэглэнэ 
\authorlong{Галсанбаяр Болд} % Таны бүтэн нэр, нүүр болон хураангуй хуудсанд ашигласан. Дурын газарт бол \longname командыг хэрэглэнэ 
\addresses{ztsbold@gmail.com} % Таны хаяг, одоогоор ашиглаагүй. Өөр газарт бол \addressname командыг хэрэглэнэ
\subject{Компьютерийн ухаан} % Таны салбар, одоогоор ашиглаагүй. Дурынр газарт бол \subjectname командыг ашиглана
\keywords{Багш, журнал} % Түлхүүр үгс, одоогоор ашиглаагүй. Дурын газарт бол \keywordnames командыг хэрэглэнэ
\university{\href{http://www.must.edu.mn}{Шинжлэх Ухаан Технологийн Их Сургууль}} % Их сургуулийн нэр ба веб хаяг. Дурын газарт бол \univname командыг хэрэглэнэ 
\department{Мэдээллийн технологийн салбар} % Сургууль/тэнхмийн нэр, нүүр болон хураангуй хуудсанд ашигласан. Дурын газарт бол\deptname командыг хэрэглэнэ
\deptchair{Доктор (Ph.D) А.Алтангэрэл} % Тэнхим/эрхлэгчийн нэр, нүүр болон хураангуй хуудсанд ашигласан. Дурын газарт бол\chairname командыг хэрэглэнэ
\group{Робот техникийн баг} % Судалгааны баг/тэнхмийн нэр, нүүр хуудсанд ашигласан. Дурын газарт бол \groupname командыг хэрэглэнэ
\faculty{\href{http://www.sict.edu.mn}{Мэдээлэл, Холбооны Технологийн Сургууль}} % Салбар сургууль/факультетийн нэр, нүүр болон хураангуй хуудсанд ашигласан. Дурын газарт бол \facname командыг ашиглана

\hypersetup{pdftitle=\ttitle} % Pdf файлын гарчиг
\hypersetup{pdfauthor=\shortname} % Pdf файлын зохиогчийн нэр
\hypersetup{pdfkeywords=\keywordnames} % Pdf файлын түлхүүр үгс
\hypersetup{allcolors=black} % Pdf файлын бүх холбоос хар өнгөтэй

\usepackage{longtable}
\begin{document}

\frontmatter % Агуулгын өмнөх хуудас дугаарлалт: i, ii, iii, iv... г.м.

\pagestyle{plain} % Тезисийн загварыг дуудах хүртэлх толгой мөрийн суурь загвар

%-------------------------------------------------------------------------------
%	TITLE PAGE
%-------------------------------------------------------------------------------
\pagecolor{Cover}
\begin{titlepage}
\begin{center}

{\scshape\LARGE \univname\par} % Их сургуулийн нэр
{\scshape\Large \facname\par}\vspace{0.5cm} % Их сургуулийн нэр

\begin{figure}[!htbp]
\centering
\includegraphics[scale=0.2]{Figures/MUST_logo.png}
\end{figure}

\vspace{1cm}
\hfill \large{\longname} \\

\vspace{1cm}

{\huge \bfseries \ttitle\par}\vspace{0.4cm} % Тезисийн нэр

\vspace{3cm}
\textsc{\Large {\thesisname}}\\ % Тезисийн төрөл

\vfill

\large {Улаанбаатар хот} \\
 
\end{center}
\end{titlepage}
\pagecolor{white}

%-------------------------------------------------------------------------------
%	SUBTITLE PAGE
%-------------------------------------------------------------------------------

\begin{titlepage}
\begin{center}

{\scshape\LARGE \univname\par} % Их сургуулийн нэр
{\scshape\Large \facname\par}\vspace{0.5cm} % Их сургуулийн нэр

\vspace{2cm}
\hfill \large{\deptname} \\

\vspace{2cm}

{\huge \bfseries \ttitle\par}\vspace{0.4cm} % Тезисийн нэр

\vspace{2cm}

\begin{minipage}[t] {0.9\textwidth}
\begin{flushleft} 
\normalsize

Мэргэжлийн индекс: \degreeid \\
Мэргэжил: \degreename \\[2cm]

\emph{Удирдагч:} {\supname} \\% Удирдагчийн нэр
%\emph{Зөвлөгч:} {\advicename} \\ % Зөвлөгч нарын нэрс
\emph{Гүйцэтгэгч:} {\shortname} \\ % Зохиогчийн нэр

\end{flushleft}
\end{minipage}

\vfill

\large {Улаанбаатар хот} \\
{\large 2018 он}\\ % Date

\end{center}
\end{titlepage}

  % Нүүр хуудас
%-------------------------------------------------------------------------------
%	WORK PLAN & REVIEW PAGE
%-------------------------------------------------------------------------------

\begin{titlepage}

\vspace*{0.5cm}
Батлав. \deptname ын эрхлэгч: 
\begin{flushright}
\makebox[4cm]{\dotfill} /\chairname/ 
\end{flushright}

Удирдагч: 
\begin{flushright}
\makebox[4cm]{\dotfill} /\supname/
\end{flushright}

\begin{center}

\vspace*{2cm}
\textbf{{\large ТӨГСӨЛТИЙН АЖЛЫН \\ ТӨЛӨВЛӨГӨӨ}}\\[0.5cm]

\textsc{\large Сэдэв: ''\ttitle''}\\[0.5cm]

\begin{tabular}{|c|p{7cm}|c|c|}
	%\rowcolor{gray!40}
	\hline
	№ & \makebox[7cm][c]{Ажлын бүлэг, хэсгийн нэр} & Эзлэх хувь & Дуусах хугацаа \\ \hline
	1 & {Диплом-Бичиг баримт Судалгаа, шинжилгээ, зохиомжийн хэсгүүд} &  70\% & 2017-10-25 \\ \hline
	2 & {Диплом-Бичиг баримт, Програм}     & 100\%,50\% & 2017-10-30 \\ \hline
	3 & {} &  &  \\ \hline
	4 & {} &  &  \\ \hline
	5 & {} &  &  \\ \hline
\end{tabular}

\vspace{2cm}
Төлөвлөгөөг боловсруулсан оюутан: \makebox[3cm]{\dotfill} /\shortname/

\end{center}

\newpage

\begin{center}

\vspace*{2cm}
\textbf{{\large ТӨГСӨЛТИЙН АЖЛЫН ЯВЦ}}\\[0.5cm]

\begin{tabular}{|c|p{7cm}|c|c|}
	%\rowcolor{gray!40}
	\hline
	№ & \makebox[7cm][c]{Хийж гүйцэтгэсэн ажил} & Биелсэн     & Удирдагчийн \\
	  &                    & хугацаа    & гарын үсэг \\ \hline
	1 & {Удиртгал}         & 2017-10-05 &  \\ \hline
	2 & {Онолын хэсэг}     & 2017-10-30 &  \\ \hline
	3 & {Судалгааны хэсэг} & 2017-11-15 &  \\ \hline
	4 & {Төслийн хэсэг}    & 2017-12-20 &  \\ \hline
	5 & {Дүгнэлт}          & 2017-12-25 &  \\ \hline
\end{tabular}

\vspace{1cm}
Ажлын товч дүгнэлт \\[0.5cm]
\dotfill \\[0.2cm]
\dotfill \\[0.2cm]
\dotfill \\[0.2cm]
\dotfill \\[0.2cm]
\dotfill \\[0.2cm]
\dotfill \\[0.2cm]
\dotfill \\[0.5cm]
Удирдагч: \makebox[3cm]{\dotfill} /\supname/ \\

\vspace{2cm}
ЗӨВШӨӨРӨЛ \\[0.5cm]
Оюутан \shortname ын бичсэн төгсөлтийн ажлыг УШК-д хамгаалуулахаар тодорхойлов.\\[0.5cm]
Салбарын эрхлэгч: \makebox[3cm]{\dotfill} /\chairname/
\end{center}

\end{titlepage}

\newpage

\begin{titlepage}
\begin{center}

{\scshape\Large \univname\par} % Их сургуулийн нэр
{\scshape\large \facname\par}\vspace{1cm} % Их сургуулийн нэр

\textbf{{\Large ШҮҮМЖИЙН ХУУДАС}}\\[1cm]

\end{center}

\normalsize

\deptname ын төгсөх курсийн оюутан \shortname -н ''\ttitle'' сэдэвт төгсөлтийн ажлын шүүмж.

\begin{enumerate}
\item Төслөөр дэвшүүлсэн асуудал, үүнтэй холбоотой онолын материал уншиж судалсан байдал. Энэ талаар хүмүүсийн хийсэн судалгаа, түүний үр дүнг уншиж тусгасан эсэх.
\begin{center}
\dotfill \\[0.1cm]
\dotfill \\[0.1cm]
\dotfill \\[0.1cm]
\dotfill \\[0.1cm]
\dotfill \\[0.1cm]
\dotfill \\[0.1cm]
\dotfill \\[0.4cm]
\end{center}
\item Төслийн ерөнхий агуулга. Шийдсэн зүйлүүд, хүрсэн үр дүн. Өөрийн санааг гарган, харьцуулалт хийн, дүгнэж байгаа чадвар.
\begin{center}
\dotfill \\[0.1cm]
\dotfill \\[0.1cm]
\dotfill \\[0.1cm]
\dotfill \\[0.1cm]
\dotfill \\[0.1cm]
\dotfill \\[0.1cm]
\dotfill \\[0.4cm]
\end{center}
\item Эмх цэгцтэй, стандарт хангасан өөрөөр хэлбэл диплом бичих шаардлагуудыг биелүүлсэн эсэх. Төсөлд анзаарагдсан алдаанууд, зөв бичгийн болон өгүүлбэр зүйн гэх мэт /Хуудас дугаарлагдаагүй, зураг хүснэгтийн дугаар болон тайлбар байхгүй, шрифт хольсон, хувилсан зүйл ихээр оруулсан/.
\begin{center}
\dotfill \\[0.1cm]
\dotfill \\[0.1cm]
\dotfill \\[0.1cm]
\dotfill \\[0.1cm]
\dotfill \\[0.1cm]
\dotfill \\[0.1cm]
\dotfill \\[0.4cm]
\end{center}
\item Төслөөр орхигдуулсан болон дутуу болсон зүйлүүд. Цаашид анхаарах хэрэгтэй зүйлүүд.
\begin{center}
\dotfill \\[0.1cm]
\dotfill \\[0.1cm]
\dotfill \\[0.1cm]
\dotfill \\[0.1cm]
\dotfill \\[0.1cm]
\dotfill \\[0.1cm]
\dotfill \\[0.4cm]
\end{center}
\item Төслийн талаар онцолж тэмдэглэх зүйлүүд.
\begin{center}
\dotfill \\[0.1cm]
\dotfill \\[0.1cm]
\dotfill \\[0.1cm]
\dotfill \\[0.1cm]
\dotfill \\[0.1cm]
\dotfill \\[0.1cm]
\dotfill \\[0.4cm]
\end{center}
\item Ерөнхий оноо. (5 оноо)
\begin{center}
\dotfill \\[1cm]
\end{center}
\end{enumerate}
Шүүмж бичсэн: \makebox[3cm]{\dotfill} /\readname/ \\[0.5cm]
Ажлын газар: \dotfill \\[0.5cm]
Хаяг (Утас) \makebox[5cm]{\dotfill}
\end{titlepage}
 % Төлөвлөгөө, гүйцэтгэл, шүүмжийн хуудас
%-------------------------------------------------------------------------------
%	DECLARATION PAGE
%-------------------------------------------------------------------------------
\begin{declaration}
\addchaptertocentry{\authorshipname}

\noindent Миний бие \shortname, ''{\ttitle}'' сэдэвт энэ ажил нь минийх бөгөөд дараахыг нотолж байна. Үүнд:

\begin{itemize} 
\item Горилогч энэ ажлыг тус сургуулиас боловсролын зэрэг авахаар бүхэлд нь буюу голлон хийсэн болно.
\item Энэ ажлын аль нэг хэсгийг тус сургуульд эсвэл өөр байгууллагад боловсролын зэрэг, мэргэшил авахаар өмнө нь илгээсэн бол түүнийгээ тодорхой заасан болно.
\item Бусад хүмүүсийн хэвлүүлсэн ажлаас зөвлөгөө авсан бол түүнийгээ үндэслэсэн болно.
\item Бусад хүмүүсийн ажлаас ишлэл хийхдээ бол эх үүсвэрийг нь заасан болно.
\item Энэхүү ажилд ашигласан бүх эх үүсвэр нь нээлттэй эхийн лицензтэй ба лиценз дээрх заалтын дагуу ашигласан болно.
\item Миний ажилд тусалсан голлох бүх эх үүсвэрт талархаж байна.
\item Ажлыг бусадтай хамтарсан бол алийг нь бусад хүмүүс хийсэн болохыг тодорхой заасан болно.
\end{itemize}
\bigskip
 
\noindent Гарын үсэг: \rule[-0.5em]{12.7em}{0.5pt}
\bigskip

\noindent Огноо: \rule[-0.5em]{15em}{0.5pt}

\end{declaration}

\clearpage

 % Мэдэгдлийн хуудас
%-------------------------------------------------------------------------------
%	QUOTATION PAGE
%-------------------------------------------------------------------------------

\vspace*{0.2\textheight}

\noindent\enquote{\itshape Thanks to my solid academic training, today I can write hundreds of words on virtually any topic without possessing a shred of information, which is how I got a good job in journalism.}\bigbreak

\hfill Dave Barry 
\normalfont

 % Ишлэл хуудас
%-------------------------------------------------------------------------------
%	ABSTRACT PAGE
%-------------------------------------------------------------------------------

\addchaptertocentry{Хураангуй} % Хураангуйг гарчигт нэмэх

\begin{center}
{\scshape\Large \univname\par} % Их сургуулийн нэр
{\scshape\large \facname\par}\vspace{0.5cm} % Их сургуулийн нэр
{\huge\textbf{{Хураангуй}} \par}
\bigskip
{\Large{\ttitle} \par} % Тезисийн нэр
\bigskip

{\normalsize \shortname \par} % Зохиогчийн нэр
\addressname
\end{center}

\textit{\textbf{Түлхүүр үгс: \keywordnames}}
\bigskip

Төслийн ажлын тухай хураангуйг энд бичих



 % Ажлын хураангуй
%-------------------------------------------------------------------------------
%	ACKNOWLEDGEMENTS
%-------------------------------------------------------------------------------

\begin{acknowledgements}
\addchaptertocentry{\acknowledgementname}

Юуны түрүүнд энэхүү төслийг гүйцэлдүүлэх явцыг удирдан явуулсан Шинжлэх Ухаан Технологийн Их Сургуулийн харьяа Мэдээлэл, Холбооны Технологийн Сургуулийн Магистр Ж.Золжаргал-д гүнээ талархал илэрхийлж байна. Мөн тус сургуулийн Компьютерийн ухааны салбарын Профессор (PhD) А.Эрдэнэбаатар, Доктор (PhD), дэд профессор Б.Батзолбоо, Доктор (PhD) А.Хүдэр, Доктор (PhD) Ч.Цэнд-Аюуш, Доктор (PhD) Ч.Золзаяа, Доктор (PhD) Д.Энхжаргал, Доктор (PhD) Ч.Эрдэнэбат, Магистр Т.Золбоо, Магистр Д.Золбоо, Магистр Б.Сод-Од, Магистр Б.Гүндсамбуу, Магистр Б.\\Мөнхбуян, Магистр Г.Цэнд-Аюуш, Магистр Б.Намжилдорж тэргүүтэй багш, профессоруудад талархал илэрхийлж байна.

\end{acknowledgements}

 % Талархлын хуудас
%-------------------------------------------------------------------------------
%	ABBREVIATIONS
%-------------------------------------------------------------------------------

\begin{abbreviations}{ll} % Товчлолын жагсаалт оруулах (хоёр багатай хүснэгт)
\addchaptertocentry{\abbrevname}

\textbf{AI} & \textbf{A}rtifical \textbf{I}ntelligence\\
\textbf{ХОУ} & \textbf{Х}иймэл \textbf{О}юун \textbf{У}хаан\\
\textbf{АНУ} & \textbf{А}мерикийн \textbf{Н}эгдсэн \textbf{У}лс\\
\textbf{JSON} & \textbf{J}ava\textbf{S}cript \textbf{O}bject \textbf{N}otation\\
\textbf{API} & \textbf{A}pplication \textbf{P}rogramming \textbf{I}nterface\\
\textbf{CPU} & \textbf{C}entral \textbf{P}rocessing \textbf{U}nit\\
\textbf{RAM} & \textbf{R}andom \textbf{A}ccess \textbf{M}emory\\


\end{abbreviations}

 % Товчилсон үгс
%%----------------------------------------------------------------------------------
%	PHYSICAL CONSTANTS/OTHER DEFINITIONS
%---------------------------------------------------------------------------------
\addchaptertocentry{\constantsname}

\begin{constants}{lr@{${}={}$}l} % Физик тогтмолыг 3 баганатай хуүснэгтээр харуулах

% The \SI{}{} command is provided by the siunitx package, see its documentation for instructions on how to use it

    \[ Биеийн жингийн индекс & BMI & \frac{\textit{масс}_\textit{кг}}{\textit{өндөр}^2_\textit{м}} = \frac{\textit{масс}_\textit{паунд}}{\textit{өндөр}^2_\textit{инч}} \times 703 \] \\
    \[ Бодисын солилцооны түвшин(Эр) & BMR & \left(\frac{10.0m}{1\textit{кг}}+\frac{6.25h}{1\textit{см}} -\frac{5.0a}{1\textit{жил}} + 5 \right) \frac{\textit{ккал}}{\textit{өдөр}} \] \\
    \[ Бодисын солилцооны түвшин(Эм) & BMR & \left(\frac{10.0m}{1\textit{кг}} + \frac{6.25h}{1\textit{см}} - \frac{5.0a}{1\textit{жил}} - 161 \right) \frac{\textit{ккал}}{\textit{өдөр}} \] \\
    \[ Өдөрт авах илчлэг & BE & \textit{Бодисын солилцооны түвшин} \times \textit{a} \] \\
%Constant Name & $Symbol$ & $Constant Value$ with units\\

\end{constants} % Ашигласан тогтмолууд
%%-------------------------------------------------------------------------------
%	SYMBOLS
%-------------------------------------------------------------------------------

\begin{symbols}{lll} % Таних тэмтгийн жагсаалт оруулах (гурван баганатай хүснэгт)
\addchaptertocentry{\symbolsname}

$a$ & distance & \si{\meter} \\
$P$ & power & \si{\watt} (\si{\joule\per\second}) \\
%Symbol & Name & Unit \\

\addlinespace % Ром тэмдгийг Грекээс зааглах

$\omega$ & angular frequency & \si{\radian} \\

\end{symbols}

 % Ашгласан таних тэмдгүүд
%%-------------------------------------------------------------------------------
%	DEDICATION
%-------------------------------------------------------------------------------

\dedicatory{Өөрийн \ldots -д зориулав} 

 % Ажлыг хэнд зориулсан

%-------------------------------------------------------------------------------
%	LIST OF CONTENTS/FIGURES/TABLES PAGES
%-------------------------------------------------------------------------------

\tableofcontents % Гарчиг хэвлэх
\listoffigures % Зургийн жагсаалт хэвлэх
\listoftables % Хүснэгтийн жагсаалт хэвлэх

%-------------------------------------------------------------------------------
%	THESIS CONTENT - CHAPTERS
%-------------------------------------------------------------------------------

\mainmatter % Хуудасны тоон (1,2,3...) дугаарлалт эхлэнэ

\pagestyle{thesis} % Хуудасны тогойг "thesis" загвар руу буцаах
\myformat % Бүлгийн нэрийг тусгай хуудсанд хэвлэх

% Тезисийн бүлгүүдийг Chapters хавтаснаас бие даасан файл байдлаар оруулах

% Бүлэг 1


%-------------------------------------------------------------------------------

% Агуулгад ашигласан хэвшүүлэлтийн зарим командын тодорхойлолт
\newcommand{\keyword}[1]{\textbf{#1}}
\newcommand{\tabhead}[1]{\textbf{#1}}
\newcommand{\code}[1]{\texttt{#1}}
\newcommand{\file}[1]{\texttt{\bfseries#1}}
\newcommand{\option}[1]{\texttt{\itshape#1}}

\section{Судалгаа}	
\subsection{Системийн судалгаа}
	Энэ төрлийн ситемүүд нь computer-ийн ухааны computer vision-ы салбар дээр үндэслэн хийгддэг бөгөөд хиймэл оюун, машин сургалттай хослон ажилдаг тэгсэнээр илүү нарийн үр дүнд хүч чаддаг. Computer vision чиглэлийн OpenCV Фреймворк нь 1999 оноос албан ёсоор нээлтээ хийж хөгжүүлэгдэж эхлэсэн хамгийн том фреймворк юм. eyeGrade програм нь хамгийн олон хэрэглэгчдэд хүрсэн desktop програм юм. Mobile төхөөрөмжийн хувьд zipgrade гэх application-ыг гадаадад ашигладаг мөн kahook гээд арай өөр төрлийн илүү өргөн хүрээг хамарсан application маш их ашиглагддаг.
\subsection{Технологийн судалгаа}
	\subsubsection{React Native}
	RN нь javascript дээр бичигдсэн ios болог android application хөгжүүлэхд зориулагдсан hybrid фреймворк юм. RN нь react javascript дээр тулгуурладаг. React - ийн component level pattern нь системийг өргөтгөхөд илүү хялбар уян хатан болгодог.
	\subsubsection{OpenCV}
	Opencv нь с болон с++  дээр бичигдсэн мөн бүх төрлийн cpu-дээр ажилладаг, бараг бүх түгээмэл програмчлалын хэл дээр байдаг. Cannyedge болон otsi-ийн илрүүлэлтийн алгоритмуудыг зэрэг маш олон хэрэгцээ ихтэй сангуудыг агуулдаг
	\subsubsection{MongoDB}
	MongoDB өгөгдөлийн сан нь харьцаа өгөдөлийн сан биш бөгөөд өгөдөлийн давхардалт зэрэг зүйлийг шалгадаггүй. Давуу тал нь илүү хурдан тооцоолох чадвартай. Бүх өгөгдөл collection-д хадгалагддаг.
	\subsubsection{Python}
	Python нь бүтээмж өндөртэй програмчлалын хэл юм.
	\subsection{Бүлгийн дүгнэлт}
	Python нь OpenCV-тэй сайн зохицож ажилладаг бөгөөд бичиглэл ойлгомжтой гэдэг үйднээс сонгосон. MongoDB бол харьцаа өгөгдөлийн сангаас хамаагүй уян хатан өөрчлөгдөх боломж хялбар бөгөөд өргөтгөхөд хялбар энгийн учир сонгосон. React-native ч мөн адил өргөтгөх өөрчлөх илүү хялбар гэдэг үүднээс сонгосон.

% Chapter Template

\section{Хөгжүүлэлтэд ашиглах технологи, хэрэгслүүд} % Main chapter title

\label{Chapter2} % Change X to a consecutive number; for referencing this chapter elsewhere, use \ref{ChapterX}

%----------------------------------------------------------------------------------------
%   SECTION 1
%----------------------------------------------------------------------------------------
\section{Технологийн тухай}
    \subsection{Технологи, прэймворк болон АПИ -ийн тухай үндсэн ойлголт}
    \subsection{Тэдгээрийн давуу ба дутагдалтай талууд}
    \subsection{Бусад ижил төстэй технологи, прэймворк, АПИ ба тэдгээртэй харьцуулсан судалгаа}
    \subsection{Технологийн сонголт дээрх SWOT шинжилгээ}
\section{Технологийн ерөнхий шаардлага}
    \subsection{Үйлдлийн системийн шаардлага}
    \subsection{Програмчлалын хэлний шаардлага}
    \subsection{Суулгах програмчлалын хэрэгслүүд}
    \subsection{Өгөгдлийн сангийн шаардлага}
\section{Суурилуулах заавар}
    \subsection{Сервер талд суурилуулах заавар}
    \subsection{Клиент талд суурилуулах заавар}
    \subsection{Програмчлалын хэрэгсэлтэй холбох заавар}
\section{Функц үйлдлүүд, тэдгээрийг хэрэглэх тухай}
%Класс, функц, массив, өгөгдлийн сантай холбох, өгөгдлийн сангаас өгөгдөл авах, бичих, устгах ... гэх мэтийг нэг бүрчлэн хэрхэн хэрэглэж байгааг тайлбарлах, жишээ гаргах
    \subsection{...}
    \subsection{...} 
% Бүлэг 3

\section{\ttitle-н шинжилгээ} % Зарим нэг зөвлөмж

\label{Chapter3} % Энэ бүлэг рүү ишлэл хийх бол \ref{Chapter3} командыг ашигла 
%-------------------------------------------------------------------------------
%	SECTION
%-------------------------------------------------------------------------------

\section{Системийн үйл ажиллагааны тухай}

%	SECTION
%-------------------------------------------------------------------------------
\section{Системийг ашиглах хэрэглэгчид}
  \subsection{Системийн оролцогч}
  \subsection{Системийн тоглогч}
  
%	SECTION
%-------------------------------------------------------------------------------
\section{Функцийн шаардлага}

%	SECTION
%-------------------------------------------------------------------------------
\section{Функцийн бус шаардлага}

%	SECTION
%-------------------------------------------------------------------------------  
\section{Юзкейс диаграм}

%	SECTION
%-------------------------------------------------------------------------------  
\section{Юзкейс тодорхойлолт}

%	SECTION
%-------------------------------------------------------------------------------
\section{Шинжилгээний класс диаграм}

%	SECTION
%-------------------------------------------------------------------------------
\section{Шинжилгээний дарааллын диаграм}

%	SECTION
%-------------------------------------------------------------------------------
\section{Үйл ажиллагааны диаграм}

%	SECTION
%-------------------------------------------------------------------------------
\section{Бүлгийн дүгнэлт}
%% Chapter Template

\section{\ttitle-н зохиомж} % Main chapter title

\label{Chapter4} % Change X to a consecutive number; for referencing this chapter elsewhere, use \ref{ChapterX}

%----------------------------------------------------------------------------------------
%   SECTION 1
%----------------------------------------------------------------------------------------
\section{Өгөгдлийн ерөнхий схем}

%!!!!!!!!!!!!!!!!!!!!!!!!!!!!!
%!!!!!!!!!!!!!!!!!!!!!!!!!!!!!
%Доорх хэсгүүдийг дипломын ажил дээр гүйцээнэ.
%!!!!!!!!!!!!!!!!!!!!!!!!!!!!!
%!!!!!!!!!!!!!!!!!!!!!!!!!!!!!
%!!!!!!!!!!!!!!!!!!!!!!!!!!!!!

%	SECTION
%-------------------------------------------------------------------------------
%\section{Класс диаграм}

%	SECTION
%-------------------------------------------------------------------------------
%\section{Дарааллын диаграм}
%Зохиомжийн хэсгийн дарааллын диаграмуудыг энд оруулна.

%\section{Төлвийн диаграм}
%Төлвийн диаграмуудыг энд оруулна.

%	SECTION
%------------------------------------------------------------------------------
%\section{Хэрэглэгчийн интерфейс}
%Системийн төсөөлж буй хэрэглэгчийн интерфейсүүдийг энд оруулна.
%	SECTION
%-------------------------------------------------------------------------------
%\section{Тайлан, статистик үзүүлэлтийн загвар}
%Системээс гарч болох тайлан эсвэл статистик үзүүлэлтийн төсөөллийг энд оруулна. Зураг байж болно.

%	SECTION
%-------------------------------------------------------------------------------
\section{Бүлгийн дүгнэлт}
 
%\include{Chapters/Chapter5}

%-------------------------------------------------------------------------------
%	THESIS CONTENT - APPENDICES
%-------------------------------------------------------------------------------

\appendix % Дараах "chapters" нь Хавсралт болохыг LaTex -д хэлэх

% Тезисийн бүлгүүдийг Appendices хавтаснаас бие даасан файл байдлаар оруулах

% Хавсралт A

%\chapter{Хавсралтын нэр} % Main appendix title

\label{AppendixA} % For referencing this appendix elsewhere, use \ref{AppendixA}

Хавсралтаа энд бичнэ.
%\include{Appendices/AppendixB}
%\include{Appendices/AppendixC}

%-------------------------------------------------------------------------------
%	BIBLIOGRAPHY
%-------------------------------------------------------------------------------
\defformat % Бүлгийн нэрийг оргиналь байдлаар хэвлэх

\addchaptertocentry{Ном зүй} % Ном зүйг гарчигт нэмэх

\printbibliography[heading=bibliography,title={Ном зүй}]

%-------------------------------------------------------------------------------
\end{document}
 